
\documentclass{article}
\usepackage{colortbl}
\usepackage{makecell}
\usepackage{multirow}
\usepackage{supertabular}

\begin{document}

\newcounter{utterance}

\centering \large Interaction Transcript for game `codenames', experiment `unambiguous', episode 0 with reflexion{-}Qwen3.5{-}27B.
\vspace{24pt}

{ \footnotesize  \setcounter{utterance}{1}
\setlength{\tabcolsep}{0pt}
\begin{supertabular}{c@{$\;$}|p{.15\linewidth}@{}p{.15\linewidth}p{.15\linewidth}p{.15\linewidth}p{.15\linewidth}p{.15\linewidth}}
    \# & $\;$A & \multicolumn{4}{c}{Game Master} & $\;\:$B\\
    \hline

    \theutterance \stepcounter{utterance}  
    & & & \multicolumn{2}{p{0.3\linewidth}}{
        \cellcolor[rgb]{0.9,0.9,0.9}{
            \makecell[{{p{\linewidth}}}]{
                \texttt{\tiny{[GM$|$GM]}}
                {'hidden': {'team': ['riddance'], 'innocent': ['underwear', 'spider', 'tornado', 'opponent', 'greed'], 'opponent': [], 'assassin': ['adultery']}, 'revealed': {'team': {'team': [], 'innocent': [], 'opponent': [], 'assassin': []}, 'opponent': {'team': [], 'innocent': [], 'opponent': [], 'assassin': []}}}
            }
        }
    }
    & & \\ \\

    \theutterance \stepcounter{utterance}  
    & & \multicolumn{4}{p{0.6\linewidth}}{
        \cellcolor[rgb]{0.9,0.9,0.9}{
            \makecell[{{p{\linewidth}}}]{
                \texttt{\tiny{[P1$\langle$GM]}}
                \texttt{You have attempted this same task before and did not succeed. These are the plans you wrote after those attempts.} \\
\texttt{Trial \#1: New plan: My previous strategy failed because I prioritized grouping words into multi{-}word clues (e.g., "suit" for blazer/councilman, "union" for merger/councilman) which inadvertently triggered guesses of distractor words ("broker") or opponent words ("lobotomy"), ending my turns prematurely. I also used overly broad clues like "water" and "death" that had high collision risks with the opponent's list. Next time, I will prioritize safety over speed by using single{-}word clues that are highly specific to one target word to minimize the chance of triggering a wrong guess. I will strictly avoid clues that could semantically link to any remaining opponent or distractor words (e.g., avoiding "suit" when "broker" is a distractor). For the final word "riddance," I will use a more precise synonym like "relief" or "freedom" rather than a generic verb like "get" to ensure the partner doesn't guess a distractor. The goal is to secure one correct guess per turn rather than risking a turn{-}ending error for two guesses.} \\
\\ 
\texttt{Trial \#2: New plan: My previous attempt failed because I provided a clue ("relief") that, while semantically correct for "riddance," was too broad and likely led my partner to guess a distractor word (such as "spider" or "tornado" via weak associations) or an opponent word, ending the turn prematurely. The strategy of prioritizing speed with broad synonyms proved too risky given the specific distractor list. Next time, I will prioritize safety by selecting a clue that is a direct, unambiguous synonym or a specific context for the target word that has zero semantic overlap with the distractors (underwear, spider, tornado, opponent, greed) or the kill word (adultery). For "riddance," instead of the generic "relief," I will use a more specific term like "freedom" or "escape" if they are safer, or strictly verify that the chosen clue cannot be associated with "spider" (e.g., avoiding "web" or "trap") or "tornado" (e.g., avoiding "storm" or "wind"). I will ensure the clue points exclusively to the target to prevent the partner from guessing a wrong word and losing the turn.} \\
\\ 
\\ 
\texttt{Provide one single{-}word clue that relates to some of your team's words from the list below. You can choose to how many and to which words your clue relates to.} \\
\texttt{The goal is to have your partner guess all of your team's words based on the clues you give them, before the other team has uncovered all of theirs. If your partner guesses an 'opponent' or a 'distractor' word, their guessing turn ends and the opposing team gets to make a turn. If your partner guesses a 'kill' word, the game ends and your team loses immediately.} \\
\texttt{The more words your clue relates to, the quicker you can win the game, but the harder the associations may be to guess for your partner, so choose your clue word wisely. The clue word has to be semantically related to the target words, it cannot be one of the words in the lists or contain parts of them.} \\
\\ 
\texttt{Always give your single{-}word clue and your comma{-}separated list of related target words in the following format and make your answers as short as possible, never include any other text than is required in this form:} \\
\\ 
\texttt{CLUE: <WORD>} \\
\texttt{TARGETS: <TARGETS>} \\
\\ 
\texttt{Your team words are: riddance.} \\
\texttt{Your opponent's team words are: .} \\
\texttt{Distractor words are: underwear, spider, tornado, opponent, greed.} \\
\texttt{The kill words are: adultery.} \\
            }
        }
    }
    & \\ \\

    \theutterance \stepcounter{utterance}  
    & \multicolumn{4}{p{0.6\linewidth}}{
        \cellcolor[rgb]{0.9,0.9,0.9}{
            \makecell[{{p{\linewidth}}}]{
                \texttt{\tiny{[P1$\rangle$GM]}}
                \texttt{CLUE: relief} \\
\texttt{TARGETS: riddance} \\
            }
        }
    }
    & & \\ \\

    \theutterance \stepcounter{utterance}  
    & & & \multicolumn{2}{p{0.3\linewidth}}{
        \cellcolor[rgb]{0.9,0.9,0.9}{
            \makecell[{{p{\linewidth}}}]{
                \texttt{\tiny{[GM$|$GM]}}
                \texttt{relief} \\
            }
        }
    }
    & & \\ \\

    \theutterance \stepcounter{utterance}  
    & & & \multicolumn{2}{p{0.3\linewidth}}{
        \cellcolor[rgb]{0.9,0.9,0.9}{
            \makecell[{{p{\linewidth}}}]{
                \texttt{\tiny{[GM$|$GM]}}
                ['riddance']
            }
        }
    }
    & & \\ \\

    \theutterance \stepcounter{utterance}  
    & & & \multicolumn{2}{p{0.3\linewidth}}{
        \cellcolor[rgb]{0.9,0.9,0.9}{
            \makecell[{{p{\linewidth}}}]{
                \texttt{\tiny{[GM$|$GM]}}
                {'word': 'riddance', 'assignment': 'team'}
            }
        }
    }
    & & \\ \\

    \theutterance \stepcounter{utterance}  
    & & & \multicolumn{2}{p{0.3\linewidth}}{
        \cellcolor[rgb]{0.9,0.9,0.9}{
            \makecell[{{p{\linewidth}}}]{
                \texttt{\tiny{[GM$|$GM]}}
                \texttt{opponent has won} \\
            }
        }
    }
    & & \\ \\

    \theutterance \stepcounter{utterance}  
    & \multicolumn{4}{p{0.6\linewidth}}{
        \cellcolor[rgb]{0.95,0.95,0.95}{
            \makecell[{{p{\linewidth}}}]{
                \texttt{\tiny{[P2$\rightarrow$P1]}}
                \texttt{Post{-}episode Reflexion (not sent to GM):} \\
\texttt{New plan: My previous attempt failed because the clue "relief," while a valid synonym for "riddance," was too broad and likely triggered a guess of a distractor word (like "spider" or "tornado") or an opponent word, ending the turn prematurely. The strategy of using generic synonyms proved too risky given the specific distractor list. Next time, I will prioritize safety by selecting a clue that is a direct, unambiguous synonym or a specific context for the target word that has zero semantic overlap with the distractors (underwear, spider, tornado, opponent, greed) or the kill word (adultery). For "riddance," instead of the generic "relief," I will use a more specific term like "freedom" or "escape" if they are safer, or strictly verify that the chosen clue cannot be associated with "spider" (e.g., avoiding "web" or "trap") or "tornado" (e.g., avoiding "storm" or "wind"). I will ensure the clue points exclusively to the target to prevent the partner from guessing a wrong word and losing the turn.} \\
            }
        }
    }
    & & \\ \\

\end{supertabular}
}

\end{document}
